\section{Actual ample original lines give whole section and metric realizations on one native model}
\label{sec:native-ample-model-section-blocks}

\begin{theorem}[Actual original positive-power geometry retains every whole section and true norm]
\label{thm:native-ample-model-section-block-map}
\lean{EnvelopeProperty.NativeAmpleModelSectionBlocks.sectionBlockEquiv,EnvelopeProperty.NativeAmpleModelSectionBlocks.sectionBlockEquiv_apply,EnvelopeProperty.NativeAmpleModelSectionBlocks.sectionBlockNorm}
\leanok
\uses{lem:positive-power-projective-block,thm:actual-module-iso-global-sections,thm:native-pulled-projective-linear-sections,thm:scheme-global-section-model-norm-isometry,thm:native-pulled-projective-section-norms}
\source{fhx-envelope-20261028:eq:FS-function}
Suppose actual original positive-power geometric data identify the first
pulled projective twist with the actual $r$th power of an original line
$L$. For every nonnegative integer $d$, the entire original global section
module of $L^{rd}$ is linearly equivalent, over the literal original field,
to the whole same fixed native model degree-$d$ section space. Every actual
original global section retains its true proper-model norm as genuine model
section norm under the same entire native analytic homeomorphism.
\end{theorem}
\begin{proof}
The genuine original exponent-block twist isomorphism identifies the whole
original $L^{rd}$ with the actual whole original pulled $O(d)$ line. Apply
its original entire global-section isomorphism, retaining the original
field action. Compose with the proved entire linear original/model section
equivalence. Its literal original coefficient formula gives the whole
forward map. The actual module isomorphism is a true global-section model
norm isometry, and the genuine whole pulled-section norm theorem identifies
the resulting actual norm with model section norm. Thus every original
section, including zero sections, has the asserted true norm equality at
every original native point. No whole-section lift or metric identity is
assumed for the model.
\end{proof}

\begin{theorem}[Every actual ample original line internally produces one whole-section metric model]
\label{thm:native-ample-model-section-blocks}
\lean{EnvelopeProperty.NativeAmpleModelSectionBlocks.exists_model_section_blocks}
\leanok
\uses{thm:native-ample-model-section-block-map,thm:native-pulled-projective-line-pivot}
\source{bj22-trivval:sec:tc}
For any actual original integral projective scheme and any actual original
geometrically ample line $L$, there exist an original positive integer $r$,
a genuine closed projective embedding of that same original scheme with the
literal original structure triangle, and a nonempty original pivot chart,
such that one same genuine model over the entire original native function
field realizes, for every $d\ge0$, the entire original global section module
of $L^{rd}$ with its original field action and every true original model norm.
There is no embedding, local section realization, finite-family lift, norm
comparison or numerical-class compatibility certificate among the inputs.
\end{theorem}
\begin{proof}
Use the actual original ample line's geometric positive-power criterion.
It internally gives $r>0$ and a true original closed projective embedding
whose actual first pulled twist is the whole $L^r$, with the original
structure triangle. The original generic point lies in some actual
projective coordinate chart, yielding an original nonempty pivot. Apply
the preceding genuine whole-section metric construction for every $d$.
All resulting sections and norms live in the same model over the entire
original native function field. This theorem concerns all actual original
exponent blocks; numerical-class independence and complete psh invariance
under changing the line representative remain further steps.
\end{proof}
